% !TEX program = xelatex
% AI-effects 프로젝트 --- AI 점유율 지표 A~C(paper 27)와 노출도 지표 3개의 직업별 격차와 격차가 큰 직업.
% 표 본문: results/table/29_reg.tex, 29_size.tex, 29_agree.tex, 29_abc.tex, 29_top_{OE,FEL,ELO}.tex,
%          29_robust_{pos,neg}.tex, 29_ksco1.tex (code/04_analysis_kor_share_exposure_gap.do가 생성)
% 전체 자료: results/table/29_kor_share_exposure_gap.xlsx
% 관련: paper 27(점유율 지표), paper 28(점유율과 노출도의 상관, 노출도 기반 점유율)
% 컴파일: xelatex (kotex 필요)
\documentclass[11pt]{article}

\usepackage{fontspec}
\usepackage{kotex}
\usepackage[a4paper,margin=2.0cm]{geometry}
\usepackage{longtable}
\usepackage{booktabs}
\usepackage{array}
\usepackage{amsmath}
\usepackage{xcolor}
\usepackage{hyperref}
\usepackage{enumitem}

\hypersetup{colorlinks=true,linkcolor=blue!50!black,urlcolor=blue!50!black}
\newcolumntype{L}[1]{>{\raggedright\arraybackslash}p{#1}}
\newcolumntype{R}[1]{>{\raggedleft\arraybackslash}p{#1}}
\setlength{\LTleft}{0pt}
\setlength{\LTright}{0pt}
\renewcommand{\arraystretch}{1.15}
\setlength{\tabcolsep}{3.5pt}

\title{AI 점유율 지표와 노출도 지표의 직업별 격차\\[2mm]
\large --- 순위 격차, 노출도 점유율과의 차이, 회귀 잔차 ---}
\author{AI-effects 프로젝트}
\date{\today}

\begin{document}
\maketitle

\begin{abstract}
\noindent paper 27의 AI 점유율 지표 A$\sim$C와 노출도 지표 3개(Observed exposure, Felten LM AIOE, Eloundou GPT-4 $\beta$)의 직업별 격차를 KSCO 8차 소분류 129개에서 구하고, 격차가 큰 직업을 찾았다. 점유율은 합이 100\%인 규모 지표이고 노출도는 1인당 강도 지표라서 그대로 뺄 수 없다. 그래서 단위를 맞추는 세 가지 방법을 썼다. (1) 순위 격차, (3) 노출도 기반 점유율과의 차이(\%p)와 로그 비율, (4) 로그 점유율을 노출도와 로그 취업자 수에 회귀한 잔차다. 모든 격차는 양수일 때 ``실제 사용이 노출도로 기대되는 것보다 크다''는 뜻이다.
\begin{itemize}[leftmargin=*]
  \item \textbf{방법마다 답하는 질문이 다르다.} 순위 격차와 회귀 잔차는 직업 규모와 거의 무관하다(취업자 수와의 순위상관 $-0.02\sim0.14$). 두 방법은 서로 잘 맞는다(0.79$\sim$0.90). 점유율 차이는 규모와 강하게 얽혀 있다($-0.57$, $-0.70$). 그래서 ``AI 영향이 노출도 기준보다 어느 직업에 몇 \%p 더 몰리는가''라는 다른 질문에 답한다. 회귀 잔차와의 순위상관도 0.42$\sim$0.47에 그친다.
  \item \textbf{A$\sim$C 중 무엇을 써도 거의 같다.} 같은 방법$\cdot$노출도에서 B와 C의 격차 순위상관은 0.99$\sim$1.00, A와 C는 0.92$\sim$0.96이다.
  \item \textbf{노출도가 같다면 직업 규모가 커도 점유율이 커지지 않는다.} 로그 점유율 회귀에서 로그 취업자 수의 계수는 $-0.14\sim0.14$로 0과 구별되지 않는다. 노출도와 규모가 설명하는 부분은 로그 점유율 분산의 35$\sim$44\%다.
  \item \textbf{여러 방법에서 공통으로 ``사용 > 노출도''인 직업}: 293 연극$\cdot$영화 및 영상 전문가(지표 C의 목록 8개 중 8개), 213 생명 및 자연과학 관련 시험원(7개), 852 냉난방 관련 설비 조작원(6개), 211 생명 및 자연과학 관련 전문가, 799 기타 기능 관련 종사자(각 5개).
  \item \textbf{공통으로 ``노출도 > 사용''인 직업}: 311 기획 및 마케팅 사무원, 531 통신 관련 판매 종사자(각 6개), 252 시민사회 활동가, 243 약사 및 한약사(각 5개).
  \item \textbf{점유율 차이로 보면} 노출도 기준보다 AI 사용이 전문가 대분류에 36$\sim$42\%p 더 몰리고, 사무 종사자에서는 21$\sim$25\%p 덜 몰린다. 상위는 291 작가 및 언론, 222 컴퓨터 시스템 및 소프트웨어, 225 정보 시스템 및 웹 운영자, 하위는 331 회계 및 경리, 311 기획 및 마케팅, 521 매장 판매다.
\end{itemize}
\end{abstract}

\section{왜 그대로 빼면 안 되는가}

\begin{itemize}[leftmargin=*]
  \item \textbf{지표 A$\sim$C}(paper 27)는 한국 전체 AI 사용 중 그 직업의 몫(\%)이다. 129개 소분류를 합하면 100이고, 직업 규모가 값에 들어 있으며, 값이 크게 쏠려 있다(상위 10개 합 약 55\%). 지표 C는 $\mathrm{FS}_j \propto \text{활용도 2}_j \times E_j$다($E_j$ = 2022년 상반기 취업자 수).
  \item \textbf{노출도}는 직업 규모와 무관한 1인당 강도다. Observed exposure는 0$\sim$0.46, Eloundou $\beta$는 0.01$\sim$0.92, Felten은 평균 0인 표준화 점수($-1.4\sim1.9$)라 단위가 서로 다르다.
  \item 그래서 $\mathrm{FS}_j - x_j$는 단위가 맞지 않는 뺄셈이다. 격차가 큰 직업이 ``취업자가 많은 직업''으로 채워진다.
\end{itemize}

\section{세 가지 방법}

소분류 $j$, 점유율 지표 $m_j \in \{\text{A, B, C}\}$(\%), 노출도 $x_j$, 순위 $r(\cdot)$(129개 안, 1 = 최상위)에 대해 다음과 같이 정의한다.

\paragraph{방법 1: 순위 격차.}
\[
  \mathrm{RG}_j = r(x_j) - r(m_j).
\]
노출도 순위보다 사용 순위가 몇 계단 높은지를 잰다. 단위와 분포 모양에 영향받지 않아 노출도 3개에 똑같이 쓸 수 있고, 값은 $-128\sim128$이다. 대신 값의 크기를 버린다. 점유율이 0.01\%와 0.02\%인 직업 사이도 한 계단이고, 11\%와 3\%인 직업 사이도 한 계단이다. 그래서 분포의 아래쪽(점유율$\cdot$노출도가 모두 작은 직업)에서 순위가 크게 움직이기 쉽다. Observed exposure는 13개 소분류가 0이라, 이들 사이의 순위(117$\sim$129위)는 정렬 순서로 정해진다. 이 13개의 RG에는 최대 12계단의 임의 부분이 있다.

\paragraph{방법 3: 노출도 기반 점유율과의 차이.} paper 28의 노출도 기반 점유율 $\mathrm{XS}^x_j = 100\,E_j x_j / \sum_k E_k x_k$를 쓴다.
\[
  \mathrm{DS}_j = m_j - \mathrm{XS}^x_j \ (\%\text{p}), \qquad \mathrm{LR}_j = \ln\!\left(m_j / \mathrm{XS}^x_j\right).
\]
두 점유율은 합이 모두 100이라 DS의 합은 0이다. DS는 ``AI 영향을 노출도로 배분했을 때와 실제 사용으로 배분했을 때, 전체의 몇 \%p가 이 직업으로 옮겨 오는가''를 뜻한다. paper 27의 시나리오 인원(점유율 $\times X$)으로 말하면, 일자리 영향 $X$명 중 노출도 기준보다 DS$\times X/100$명이 더(또는 덜) 이 직업에 몰린다는 뜻이다. 다만 규모가 큰 직업일수록 차이가 커지므로 상위$\cdot$하위가 큰 직업으로 채워진다.

LR은 비율이라 규모가 약분된다. 지표 C에서는 $\mathrm{LR}_j = \ln(\text{활용도 2}_j / x_j) + \text{상수}$가 된다(A는 활용도 1, B는 1인당 절감 시간). 즉 강도끼리의 비율로 돌아가므로, 취업자가 적은 직업에서 값이 커지는 문제(paper 27이 점유율 지표를 만든 이유)가 다시 생긴다. Observed exposure가 0인 13개는 정의되지 않는다. Felten은 음수가 있어 XS를 만들 수 없으므로 방법 3은 Observed exposure와 Eloundou $\beta$에만 적용했다. 이 노트에서 LR은 참고용이고(상관표와 엑셀), 상위$\cdot$하위 목록은 DS로 만들었다.

\paragraph{방법 4: 회귀 잔차.}
\[
  \ln m_j = a + b\,x_j + c\,\ln E_j + \varepsilon_j, \qquad \mathrm{RE}_j = \hat{\varepsilon}_j.
\]
점유율 지표 3개 $\times$ 노출도 3개, 모두 9개 회귀를 OLS로 추정했다(표준오차는 robust). 잔차는 ``노출도와 직업 규모로 예측한 점유율보다 실제 점유율이 몇 배인가''를 로그로 나타낸다. 예를 들어 RE = 1이면 예측치의 $e^1 \approx 2.7$배, $-1$이면 0.37배다. 규모를 $c$로 자료에서 정하므로 LR($c = 1$, 노출도 계수도 로그 1로 고정한 경우)보다 덜 제약적이다. 대신 결과가 함수형(노출도는 선형, 규모는 로그)에 따라 달라진다.

\begin{longtable}{L{2.4cm}L{3.6cm}L{2.0cm}L{2.0cm}L{6.2cm}}
\caption{세 방법의 특성}\label{tab:methods}\\
\toprule
방법 & 정의 & 노출도 & 규모 영향 & 답하는 질문 / 약점 \\
\midrule
\endfirsthead
(1) 순위 격차 & $r(x) - r(m)$ & 3개 모두 & 거의 없음 & 분포 안 위치가 얼마나 다른가 / 값의 크기를 버림, 하위권에서 크게 흔들림 \\
(3) 점유율 차이 & $m - \mathrm{XS}^x$ (\%p) & OE, $\beta$ & 큼(큰 직업이 상위$\cdot$하위) & 노출도 대신 실제 사용으로 배분하면 몇 \%p가 옮겨 가는가 / 작은 직업의 격차는 드러나지 않음 \\
(3) 로그 비율 & $\ln(m / \mathrm{XS}^x)$ & OE, $\beta$ & 작은 직업이 커짐 & 사용 강도 / 노출도 / 강도 지표의 단점이 되살아남, OE = 0이면 정의 안 됨 \\
(4) 회귀 잔차 & $\ln m$을 $x$, $\ln E$에 회귀한 잔차 & 3개 모두 & 없음(구성상) & 노출도$\cdot$규모로 예측한 것보다 몇 배 많이 쓰는가 / 함수형에 의존 \\
\bottomrule
\end{longtable}

\begin{longtable}{L{3.4cm}L{3.4cm}R{1.4cm}R{1.4cm}R{1.6cm}R{1.4cm}R{1.0cm}R{0.8cm}}
\caption{방법 4 회귀: $\ln$(점유율) $= a + b \cdot$ 노출도 $+ c \cdot \ln$(취업자) (괄호 안은 robust 표준오차)}\label{tab:reg}\\
\toprule
점유율 지표 & 노출도 & $b$ & (s.e.) & $c$ & (s.e.) & $R^2$ & $N$ \\
\midrule
\endfirsthead
% 04_12_analysis_kor_share_exposure_gap.do가 생성. 직접 고치지 말 것.
A 대화 점유율 & Observed exposure & 12.31 & (1.34) & -0.00 & (0.11) & 0.38 & 129 \\
 & Felten LM AIOE & 1.58 & (0.19) & 0.14 & (0.11) & 0.35 & 129 \\
 & Eloundou GPT-4 $\beta$ & 7.14 & (0.84) & 0.11 & (0.12) & 0.43 & 129 \\
\midrule
B 절감 시간 점유율 & Observed exposure & 14.43 & (1.56) & -0.15 & (0.14) & 0.37 & 129 \\
 & Felten LM AIOE & 1.98 & (0.24) & 0.03 & (0.14) & 0.39 & 129 \\
 & Eloundou GPT-4 $\beta$ & 8.73 & (1.05) & -0.01 & (0.14) & 0.44 & 129 \\
\midrule
C 취업자 환산 점유율 & Observed exposure & 14.35 & (1.57) & -0.14 & (0.14) & 0.36 & 129 \\
 & Felten LM AIOE & 1.98 & (0.24) & 0.04 & (0.14) & 0.39 & 129 \\
 & Eloundou GPT-4 $\beta$ & 8.63 & (1.05) & -0.00 & (0.14) & 0.43 & 129 \\

\bottomrule
\end{longtable}

\begin{itemize}[leftmargin=*]
  \item 노출도 계수는 모두 양수이고 뚜렷하다. 지표 C 기준으로 Observed exposure가 0.1 높으면 점유율이 $e^{1.40} \approx 4.1$배, Eloundou $\beta$가 0.1 높으면 $e^{0.84} \approx 2.3$배다.
  \item \textbf{로그 취업자 수의 계수는 0과 구별되지 않는다}($-0.14\sim0.14$). 노출도가 같다면 취업자가 두 배인 직업도 AI 사용 점유율이 커지지 않는다는 뜻이다. 다시 말해 큰 직업일수록 1인당 사용이 적다. 이 때문에 잔차는 사실상 $\ln m - b\,x$에 가깝고, 규모 통제의 역할은 작다.
  \item $R^2$는 0.35$\sim$0.44다. 로그 점유율 분산의 절반 넘게가 노출도와 규모로 설명되지 않는다.
\end{itemize}

\begin{longtable}{L{4.6cm}R{2.0cm}R{2.0cm}R{2.0cm}}
\caption{지표 C 격차와 취업자 수의 Spearman 순위상관}\label{tab:size}\\
\toprule
방법 & OE & Felten & $\beta$ \\
\midrule
\endfirsthead
% 04_analysis_kor_share_exposure_gap.do가 생성. 직접 고치지 말 것.
순위 격차(1) & -0.02 & 0.11 & 0.14 \\
점유율 차이(3) & -0.57 & -- & -0.70 \\
로그 비율(3) & -0.65 & -- & -0.57 \\
회귀 잔차(4) & 0.01 & 0.02 & 0.05 \\

\bottomrule
\end{longtable}

표~\ref{tab:size}는 방법별로 격차에 규모가 얼마나 남는지 보여 준다. 순위 격차와 회귀 잔차는 규모와 거의 무관하다. 점유율 차이는 큰 직업일수록 음수(노출도 > 사용)이고, 로그 비율은 작은 직업일수록 양수다.

\section{방법 간 일치}

\begin{longtable}{L{5.4cm}R{2.0cm}R{2.0cm}R{2.0cm}}
\caption{지표 C 격차의 방법 간 Spearman 순위상관}\label{tab:agree}\\
\toprule
방법 쌍 & OE & Felten & $\beta$ \\
\midrule
\endfirsthead
% 04_analysis_kor_share_exposure_gap.do가 생성. 직접 고치지 말 것.
순위 격차(1) -- 점유율 차이(3) & 0.55 & -- & 0.28 \\
순위 격차(1) -- 로그 비율(3) & 0.66 & -- & 0.37 \\
순위 격차(1) -- 회귀 잔차(4) & 0.79 & 0.90 & 0.89 \\
점유율 차이(3) -- 로그 비율(3) & 0.86 & -- & 0.81 \\
점유율 차이(3) -- 회귀 잔차(4) & 0.47 & -- & 0.42 \\
로그 비율(3) -- 회귀 잔차(4) & 0.63 & -- & 0.63 \\

\bottomrule
\end{longtable}

\begin{longtable}{L{3.4cm}R{1.8cm}R{1.8cm}R{1.8cm}R{1.8cm}R{1.8cm}R{1.8cm}}
\caption{같은 방법$\cdot$노출도에서 지표 A, B의 격차와 지표 C의 격차의 Spearman 순위상관}\label{tab:abc}\\
\toprule
 & \multicolumn{2}{c}{OE} & \multicolumn{2}{c}{Felten} & \multicolumn{2}{c}{$\beta$} \\
방법 & A--C & B--C & A--C & B--C & A--C & B--C \\
\midrule
\endfirsthead
% 04_analysis_kor_share_exposure_gap.do가 생성. 직접 고치지 말 것.
순위 격차(1) & 0.93 & 1.00 & 0.95 & 1.00 & 0.93 & 1.00 \\
점유율 차이(3) & 0.94 & 1.00 & -- & -- & 0.96 & 0.99 \\
로그 비율(3) & 0.95 & 1.00 & -- & -- & 0.96 & 1.00 \\
회귀 잔차(4) & 0.93 & 1.00 & 0.94 & 1.00 & 0.92 & 1.00 \\

\bottomrule
\end{longtable}

\begin{itemize}[leftmargin=*]
  \item \textbf{순위 격차와 회귀 잔차는 같은 것을 잰다}(0.79$\sim$0.90). 둘 다 규모와 무관하게 ``노출도로 본 위치 대비 사용의 위치''를 본다. 회귀 잔차는 크기 정보를 쓰고, 순위 격차는 순서만 쓴다는 점이 다르다.
  \item \textbf{점유율 차이는 다른 질문에 답한다.} 순위 격차, 회귀 잔차와의 상관이 0.28$\sim$0.55로 낮다. 점유율 차이는 로그 비율과 0.81$\sim$0.86으로 맞는데, 두 방법 모두 노출도 기반 점유율을 기준으로 삼기 때문이다.
  \item \textbf{A$\sim$C의 선택은 결과를 거의 바꾸지 않는다.} B와 C는 0.99$\sim$1.00, A와 C는 0.92$\sim$0.96이다. 대화 수(A)와 절감 시간(B, C)의 차이는 대화당 절감 시간의 직업 차이뿐이고, 노출도와의 격차에 비하면 작다. 아래 목록은 지표 C로 만들었고, A와 B의 목록은 엑셀 시트 \texttt{top10}에 있다.
\end{itemize}

\section{격차가 큰 직업}

표~\ref{tab:topOE}--\ref{tab:topELO}는 지표 C의 격차 상위 10개(사용 > 노출도)와 하위 10개(노출도 > 사용)다. 괄호 안 값은 격차다(순위 격차는 계단, 점유율 차이는 \%p, 회귀 잔차는 로그).

{\small
\begin{longtable}{R{0.5cm}L{5.4cm}L{5.4cm}L{5.4cm}}
\caption{지표 C와 Observed exposure의 격차 상위$\cdot$하위 10}\label{tab:topOE}\\
\toprule
 & 순위 격차(1) & 점유율 차이(3, \%p) & 회귀 잔차(4) \\
\midrule
\endfirsthead
\toprule
 & 순위 격차(1) & 점유율 차이(3, \%p) & 회귀 잔차(4) \\
\midrule
\endhead
% 04_analysis_kor_share_exposure_gap.do가 생성. 직접 고치지 말 것.
\multicolumn{4}{l}{\textit{사용 > 노출도(격차 상위 10)}} \\
1 & 852 냉난방 관련 설비 조작원 (+94) & 291 작가 및 언론 관련 전문가 (+9.85) & 293 연극‧영화 및 영상 전문가 (+3.63) \\
2 & 149 기타 건설‧전기 및 생산 관련 관리자 (+77) & 222 컴퓨터 시스템 및 소프트웨어 전문가 (+4.09) & 235 기계‧로봇공학 기술자 및 시험원 (+3.52) \\
3 & 237 환경공학‧가스‧에너지 기술자 및 시험원 (+60) & 283 상품 기획‧홍보 및 조사 전문가 (+4.04) & 264 문리‧기술 및 예능 강사 (+3.20) \\
4 & 235 기계‧로봇공학 기술자 및 시험원 (+57) & 332 통계‧데이터 관련 사무원 (+3.97) & 239 기타 공학 전문가 및 관련 종사자 (+3.18) \\
5 & 239 기타 공학 전문가 및 관련 종사자 (+49) & 225 정보 시스템 및 웹 운영자 (+3.84) & 211 생명 및 자연과학 관련 전문가 (+3.13) \\
6 & 293 연극‧영화 및 영상 전문가 (+47) & 213 생명 및 자연과학 관련 시험원 (+3.44) & 852 냉난방 관련 설비 조작원 (+2.82) \\
7 & 236 소방‧방재‧산업안전 기술자 및 시험원 (+45) & 223 네트워크 및 정보 보안 전문가 (+2.90) & 131 연구‧교육 및 법률 관련 관리자 (+2.81) \\
8 & 246 치료‧재활사 및 의료기사 (+45) & 293 연극‧영화 및 영상 전문가 (+2.53) & 236 소방‧방재‧산업안전 기술자 및 시험원 (+2.80) \\
9 & 861 발전 및 배전 장치 조작원 (+42) & 211 생명 및 자연과학 관련 전문가 (+2.44) & 213 생명 및 자연과학 관련 시험원 (+2.67) \\
10 & 211 생명 및 자연과학 관련 전문가 (+38) & 292 학예사‧사서 및 기록물 관리사 (+2.11) & 799 기타 기능 관련 종사자 (+2.54) \\
\midrule
\multicolumn{4}{l}{\textit{노출도 > 사용(격차 하위 10)}} \\
1 & 532 방문 및 노점 판매 종사자 (-73) & 521 매장 판매 종사자 (-8.23) & 451 조리사 (-5.95) \\
2 & 441 운송 서비스 종사자 (-66) & 331 회계 및 경리 사무원 (-8.06) & 620 임업 관련 종사자 (-5.11) \\
3 & 322 운송 및 무역 사무원 (-65) & 371 총무 사무원 및 대학 행정조교 (-6.33) & 952 음식 관련 단순 종사자 (-4.78) \\
4 & 243 약사 및 한약사 (-61) & 311 기획 및 마케팅 사무원 (-5.76) & 630 어업 관련 종사자 (-3.79) \\
5 & 245 영양사 (-61) & 262 학교 교사 (-3.62) & 922 배달원 (-3.64) \\
6 & 531 통신 관련 판매 종사자 (-60) & 510 영업 종사자 (-3.55) & 441 운송 서비스 종사자 (-3.41) \\
7 & 311 기획 및 마케팅 사무원 (-57) & 362 고객 상담 및 모니터 요원 (-3.51) & 871 철도‧전동차 기관사 및 관련 종사자 (-3.24) \\
8 & 362 고객 상담 및 모니터 요원 (-49) & 284 감정‧기술 영업 및 중개 관련 종사자 (-3.24) & 252 시민사회 활동가 (-3.20) \\
9 & 244 간호사 (-47) & 531 통신 관련 판매 종사자 (-3.06) & 611 작물 재배 종사자 (-3.18) \\
10 & 922 배달원 (-47) & 312 인사 및 교육·훈련 사무원 (-2.55) & 873 물품 이동 장비 조작원 (-3.07) \\

\bottomrule
\end{longtable}

\begin{longtable}{R{0.5cm}L{5.4cm}L{5.4cm}L{5.4cm}}
\caption{지표 C와 Felten LM AIOE의 격차 상위$\cdot$하위 10(점유율 차이는 정의되지 않음)}\label{tab:topFEL}\\
\toprule
 & 순위 격차(1) & 점유율 차이(3, \%p) & 회귀 잔차(4) \\
\midrule
\endfirsthead
\toprule
 & 순위 격차(1) & 점유율 차이(3, \%p) & 회귀 잔차(4) \\
\midrule
\endhead
% 04_analysis_kor_share_exposure_gap.do가 생성. 직접 고치지 말 것.
\multicolumn{4}{l}{\textit{사용 > 노출도(격차 상위 10)}} \\
1 & 799 기타 기능 관련 종사자 (+84) & -- & 213 생명 및 자연과학 관련 시험원 (+4.08) \\
2 & 213 생명 및 자연과학 관련 시험원 (+76) & -- & 799 기타 기능 관련 종사자 (+3.87) \\
3 & 852 냉난방 관련 설비 조작원 (+61) & -- & 852 냉난방 관련 설비 조작원 (+3.32) \\
4 & 293 연극‧영화 및 영상 전문가 (+57) & -- & 222 컴퓨터 시스템 및 소프트웨어 전문가 (+3.26) \\
5 & 730 목재‧가구‧악기 및 간판 관련 기능 종사자 (+55) & -- & 225 정보 시스템 및 웹 운영자 (+3.07) \\
6 & 292 학예사‧사서 및 기록물 관리사 (+48) & -- & 293 연극‧영화 및 영상 전문가 (+3.05) \\
7 & 149 기타 건설‧전기 및 생산 관련 관리자 (+46) & -- & 372 비서 및 사무 지원 종사자 (+3.04) \\
8 & 782 건설 관련 기능 종사자 (+45) & -- & 292 학예사‧사서 및 기록물 관리사 (+2.73) \\
9 & 999 기타 서비스 관련 단순 종사자 (+44) & -- & 730 목재‧가구‧악기 및 간판 관련 기능 종사자 (+2.67) \\
10 & 225 정보 시스템 및 웹 운영자 (+43) & -- & 291 작가 및 언론 관련 전문가 (+2.64) \\
\midrule
\multicolumn{4}{l}{\textit{노출도 > 사용(격차 하위 10)}} \\
1 & 532 방문 및 노점 판매 종사자 (-94) & -- & 451 조리사 (-6.01) \\
2 & 245 영양사 (-91) & -- & 252 시민사회 활동가 (-4.61) \\
3 & 531 통신 관련 판매 종사자 (-71) & -- & 952 음식 관련 단순 종사자 (-4.33) \\
4 & 252 시민사회 활동가 (-69) & -- & 532 방문 및 노점 판매 종사자 (-3.89) \\
5 & 254 종교 관련 종사자 (-69) & -- & 245 영양사 (-3.71) \\
6 & 243 약사 및 한약사 (-61) & -- & 620 임업 관련 종사자 (-3.49) \\
7 & 322 운송 및 무역 사무원 (-59) & -- & 243 약사 및 한약사 (-3.39) \\
8 & 253 사회복지 전문가 및 관련 종사자 (-53) & -- & 531 통신 관련 판매 종사자 (-3.22) \\
9 & 263 유치원 교사 (-50) & -- & 922 배달원 (-2.95) \\
10 & 311 기획 및 마케팅 사무원 (-49) & -- & 630 어업 관련 종사자 (-2.94) \\

\bottomrule
\end{longtable}

\begin{longtable}{R{0.5cm}L{5.4cm}L{5.4cm}L{5.4cm}}
\caption{지표 C와 Eloundou GPT-4 $\beta$의 격차 상위$\cdot$하위 10}\label{tab:topELO}\\
\toprule
 & 순위 격차(1) & 점유율 차이(3, \%p) & 회귀 잔차(4) \\
\midrule
\endfirsthead
\toprule
 & 순위 격차(1) & 점유율 차이(3, \%p) & 회귀 잔차(4) \\
\midrule
\endhead
% 04_12_analysis_kor_share_exposure_gap.do가 생성. 직접 고치지 말 것.
\multicolumn{4}{l}{\textit{사용 > 노출도(격차 상위 10)}} \\
1 & 852 냉난방 관련 설비 조작원 (+78) & 291 작가 및 언론 관련 전문가 (+10.29) & 852 냉난방 관련 설비 조작원 (+3.55) \\
2 & 799 기타 기능 관련 종사자 (+74) & 222 컴퓨터 시스템 및 소프트웨어 전문가 (+7.78) & 799 기타 기능 관련 종사자 (+3.47) \\
3 & 730 목재‧가구‧악기 및 간판 관련 기능 종사자 (+56) & 283 상품 기획‧홍보 및 조사 전문가 (+4.59) & 293 연극‧영화 및 영상 전문가 (+3.11) \\
4 & 293 연극‧영화 및 영상 전문가 (+55) & 225 정보 시스템 및 웹 운영자 (+4.47) & 264 문리‧기술 및 예능 강사 (+3.00) \\
5 & 264 문리‧기술 및 예능 강사 (+49) & 332 통계‧데이터 관련 사무원 (+4.05) & 213 생명 및 자연과학 관련 시험원 (+2.86) \\
6 & 999 기타 서비스 관련 단순 종사자 (+49) & 213 생명 및 자연과학 관련 시험원 (+3.45) & 730 목재‧가구‧악기 및 간판 관련 기능 종사자 (+2.59) \\
7 & 951 가사 및 육아 도우미 (+48) & 223 네트워크 및 정보 보안 전문가 (+3.23) & 269 기타 교육 전문가 (+2.32) \\
8 & 213 생명 및 자연과학 관련 시험원 (+45) & 372 비서 및 사무 지원 종사자 (+3.21) & 510 영업 종사자 (+2.27) \\
9 & 510 영업 종사자 (+42) & 293 연극‧영화 및 영상 전문가 (+2.37) & 951 가사 및 육아 도우미 (+2.25) \\
10 & 851 금속 공작 기계 조작원 (+41) & 211 생명 및 자연과학 관련 전문가 (+2.35) & 999 기타 서비스 관련 단순 종사자 (+2.22) \\
\midrule
\multicolumn{4}{l}{\textit{노출도 > 사용(격차 하위 10)}} \\
1 & 311 기획 및 마케팅 사무원 (-82) & 331 회계 및 경리 사무원 (-7.54) & 451 조리사 (-6.39) \\
2 & 245 영양사 (-79) & 311 기획 및 마케팅 사무원 (-6.87) & 620 임업 관련 종사자 (-4.85) \\
3 & 322 운송 및 무역 사무원 (-72) & 521 매장 판매 종사자 (-5.52) & 311 기획 및 마케팅 사무원 (-4.58) \\
4 & 531 통신 관련 판매 종사자 (-64) & 321 자재 및 생산관리 사무원 (-3.85) & 252 시민사회 활동가 (-4.26) \\
5 & 252 시민사회 활동가 (-57) & 872 자동차 운전원 (-3.48) & 630 어업 관련 종사자 (-4.17) \\
6 & 243 약사 및 한약사 (-56) & 371 총무 사무원 및 대학 행정조교 (-2.75) & 952 음식 관련 단순 종사자 (-3.78) \\
7 & 331 회계 및 경리 사무원 (-53) & 352 정부 행정 사무 종사자 (-2.21) & 922 배달원 (-3.66) \\
8 & 362 고객 상담 및 모니터 요원 (-47) & 611 작물 재배 종사자 (-2.13) & 611 작물 재배 종사자 (-3.51) \\
9 & 134 문화‧예술 관련 관리자 (-46) & 531 통신 관련 판매 종사자 (-2.02) & 872 자동차 운전원 (-3.34) \\
10 & 159 기타 판매 및 고객 서비스 관리자 (-44) & 362 고객 상담 및 모니터 요원 (-1.72) & 243 약사 및 한약사 (-3.21) \\

\bottomrule
\end{longtable}
}

\begin{itemize}[leftmargin=*]
  \item \textbf{점유율 차이}: 상위는 노출도 3개에서 모두 정보통신$\cdot$글쓰기 전문가다(291 작가 및 언론 +9.9$\sim$+10.3\%p, 222 컴퓨터 시스템 및 소프트웨어, 225 정보 시스템 및 웹 운영자, 283 상품 기획$\cdot$홍보 및 조사, 332 통계$\cdot$데이터 사무원). 하위는 취업자가 많은 사무$\cdot$판매직이다(331 회계 및 경리 $-7.5\sim-8.1$\%p, 311 기획 및 마케팅, 521 매장 판매, 371 총무 사무원, 321 자재 및 생산관리 사무원). paper 28이 노출도 기반 점유율 상위로 짚은 직업들이 그대로 나온다.
  \item \textbf{순위 격차}: 상위에 노출도가 아주 낮은 기능$\cdot$장치 조작$\cdot$단순노무 직업(852 냉난방 관련 설비 조작원, 799 기타 기능 관련 종사자, 730 목재$\cdot$가구$\cdot$악기 및 간판 기능, 999 기타 서비스 관련 단순 종사자, 951 가사 및 육아 도우미)과 Observed exposure가 낮은 공학 기술자(235, 236, 237, 239)가 든다. 이들의 점유율은 0.1$\sim$0.8\%로 크지 않지만, 노출도가 최하위권이라 순위 차이가 커진다. 하위에는 노출도는 상위권인데 사용이 적은 직업(311 기획 및 마케팅, 322 운송 및 무역 사무원, 245 영양사, 243 약사, 531 통신 관련 판매, 252 시민사회 활동가, 362 고객 상담)이 든다.
  \item \textbf{회귀 잔차}: 상위는 순위 격차와 많이 겹친다(852, 799, 293, 213, 211, 264 문리$\cdot$기술 및 예능 강사). Felten 기준으로는 222, 225, 291 같은 정보통신 직업도 상위에 든다. Felten은 이들 직업의 노출도를 상대적으로 낮게 본다(paper 28). 하위는 실제 사용이 거의 없는 대인$\cdot$현장 직업(451 조리사 $-5.8\sim-6.0$, 952 음식 관련 단순 종사자, 922 배달원, 620 임업, 630 어업, 872 자동차 운전원)이다. 지표 C 점유율이 모두 0.002\% 이하다.
\end{itemize}

\section{여러 방법에서 공통으로 격차가 큰 직업}

지표 C의 목록 8개(순위 격차 3개, 점유율 차이 2개, 회귀 잔차 3개) 중 상위 10 또는 하위 10에 5번 이상 든 직업이다. 로그 비율은 강도 비율이라 세지 않았다. 순위 열은 129개 중 순위, 회귀 잔차는 Eloundou $\beta$ 기준이다.

\begin{longtable}{L{5.0cm}R{1.2cm}R{1.2cm}R{1.0cm}R{1.0cm}R{1.2cm}R{1.0cm}R{1.6cm}R{1.6cm}}
\caption{여러 방법에서 공통으로 격차가 큰 소분류(지표 C, 목록 8개 중)}\label{tab:robust}\\
\toprule
소분류 & 횟수 & C(\%) & C 순위 & OE 순위 & Felten 순위 & $\beta$ 순위 & 잔차($\beta$) & 취업자(천 명) \\
\midrule
\endfirsthead
\multicolumn{9}{l}{\textit{사용 > 노출도(상위 10에 든 횟수)}} \\
% 04_analysis_kor_share_exposure_gap.do가 생성. 직접 고치지 말 것.
293 연극‧영화 및 영상 전문가 & 8 & 2.91 & 9 & 56 & 66 & 64 & +3.08 & 120 \\
213 생명 및 자연과학 관련 시험원 & 7 & 3.51 & 7 & 40 & 83 & 52 & +2.85 & 11 \\
852 냉난방 관련 설비 조작원 & 6 & 0.63 & 35 & 129 & 96 & 113 & +3.47 & 17 \\
211 생명 및 자연과학 관련 전문가 & 5 & 2.78 & 10 & 48 & 43 & 31 & +2.19 & 76 \\
799 기타 기능 관련 종사자 & 5 & 0.78 & 28 & 62 & 112 & 102 & +3.39 & 70 \\

\midrule
\multicolumn{9}{l}{\textit{노출도 > 사용(하위 10에 든 횟수)}} \\
% 04_analysis_kor_share_exposure_gap.do가 생성. 직접 고치지 말 것.
311 기획 및 마케팅 사무원 & 6 & 0.06 & 84 & 27 & 35 & 2 & -4.52 & 707 \\
531 통신 관련 판매 종사자 & 6 & 0.11 & 72 & 12 & 1 & 8 & -2.92 & 255 \\
252 시민사회 활동가 & 5 & 0.00 & 121 & 118 & 52 & 63 & -4.27 & 4 \\
243 약사 및 한약사 & 5 & 0.01 & 111 & 50 & 50 & 55 & -3.23 & 42 \\

\bottomrule
\end{longtable}

\begin{itemize}[leftmargin=*]
  \item \textbf{293 연극$\cdot$영화 및 영상 전문가}는 8개 목록 모두에서 상위 10에 든다. 지표 C 9위(2.9\%)인데 노출도 순위는 56$\sim$66위다. 다만 paper 26에서 이 소분류 절감 시간의 46.9\%(대화량의 56.4\%)는 27-3023 News Analysts, Reporters, and Journalists에서 연계표로 배분된 것이다. 기자$\cdot$언론 과업의 사용이 연계 경로로 들어온 몫이 크다.
  \item \textbf{213 생명 및 자연과학 관련 시험원, 211 생명 및 자연과학 관련 전문가}: paper 26에서 213 절감 시간의 38.0\%는 15-2051 Data Scientists에서 연계표로 배분된 것이다. 이 몫은 생명$\cdot$자연과학 시험 과업의 사용이 아니므로 격차의 일부는 연계 경로 때문이다. 211은 19-1029 Biological Scientists, All Other(39.1\%)가 1위라 213보다는 직업 내용에 가깝다. 213은 취업자가 1만 1천 명으로 아주 작은데 점유율이 3.5\%다.
  \item \textbf{852 냉난방 관련 설비 조작원, 799 기타 기능 관련 종사자}는 노출도가 낮은데(852는 Observed exposure가 0(129위)이고 Felten$\cdot$$\beta$도 96$\cdot$113위, 799는 Felten$\cdot$$\beta$ 기준 100위 밖) 점유율이 0.6$\sim$0.8\%다. 어떤 SOC의 대화가 연계표로 배분됐는지는 아직 확인하지 않았다(paper 26 방식으로 확인 필요). 연계 경로 때문일 가능성이 있다.
  \item \textbf{311 기획 및 마케팅 사무원}은 Eloundou $\beta$ 2위인데 지표 C 84위다. paper 28에서 짚었듯이 마케팅 관련 대화가 전문가 SOC로 분류되어 283 상품 기획$\cdot$홍보 및 조사 전문가로 갔을 수 있다. 실제 사용이 적어서인지 분류 경로 때문인지 구별되지 않는다.
  \item \textbf{531 통신 관련 판매 종사자, 243 약사 및 한약사, 252 시민사회 활동가}는 노출도(특히 Felten)는 높지만 실제 사용 점유율이 0.11\% 이하다.
\end{itemize}

\section{대분류별 격차}

\begin{longtable}{L{3.4cm}R{0.7cm}R{1.0cm}R{1.0cm}R{1.0cm}R{1.0cm}R{1.2cm}R{1.2cm}R{1.0cm}R{1.0cm}R{1.0cm}}
\caption{대분류별 지표 C 격차(순위 격차$\cdot$회귀 잔차는 소분류 단순 평균, 점유율 차이는 합)}\label{tab:ksco1}\\
\toprule
 & & & \multicolumn{3}{c}{순위 격차(1)} & \multicolumn{2}{c}{점유율 차이(3, \%p)} & \multicolumn{3}{c}{회귀 잔차(4)} \\
대분류 & 소분류 수 & C(\%) & OE & Felten & $\beta$ & OE & $\beta$ & OE & Felten & $\beta$ \\
\midrule
\endfirsthead
% 04_12_analysis_kor_share_exposure_gap.do가 생성. 직접 고치지 말 것.
1 관리자 & 16 & 8.4 & +3.6 & -10.7 & -5.8 & +6.07 & +6.07 & +0.79 & -0.01 & +0.39 \\
2 전문가 및 관련 종사자 & 45 & 71.9 & +4.1 & +1.0 & +1.3 & +36.07 & +41.55 & +0.75 & +0.47 & +0.51 \\
3 사무 종사자 & 13 & 15.0 & -15.5 & -10.8 & -23.2 & -25.13 & -20.72 & +0.02 & +0.06 & -0.73 \\
4 서비스 종사자 & 13 & 0.6 & -9.0 & -13.7 & -3.0 & -1.26 & -5.80 & -1.04 & -1.30 & -0.75 \\
5 판매 종사자 & 5 & 1.2 & -35.2 & -21.4 & -8.8 & -15.37 & -9.14 & -0.71 & -0.86 & -0.03 \\
6 농림∙어업 숙련 종사자 & 4 & 0.0 & -0.5 & -4.3 & -13.0 & +0.00 & -2.49 & -4.07 & -3.15 & -3.65 \\
7 기능원 및 관련 기능 종사자 & 14 & 1.5 & +4.6 & +26.8 & +22.6 & +0.33 & -1.55 & +0.20 & +1.18 & +0.87 \\
8 장치∙기계 조작 및 조립 종사자 & 10 & 0.9 & +13.3 & +4.6 & +3.6 & +0.40 & -3.53 & -1.09 & -0.50 & -0.72 \\
9 단순노무 종사자 & 9 & 0.5 & +6.2 & +16.2 & +12.9 & -1.12 & -4.39 & -0.59 & +0.05 & -0.03 \\

\bottomrule
\end{longtable}

\begin{itemize}[leftmargin=*]
  \item \textbf{점유율 차이}: 전문가 대분류는 노출도 기준보다 36$\sim$42\%p 더 받고, 사무(약 $-21\sim-25$\%p)와 판매($-9\sim-15$\%p)는 덜 받는다. AI 영향의 배분을 노출도 대신 실제 사용으로 바꾸면 사무$\cdot$판매에서 전문가로 크게 옮겨 간다.
  \item \textbf{순위 격차}: 사무$\cdot$판매가 평균 음수, 기능원$\cdot$장치 조작$\cdot$단순노무가 평균 양수다. 다만 후자는 노출도가 최하위권이라 순위가 쉽게 올라가는 영향이 크다. 점유율로는 셋을 합쳐 2.9\%다.
  \item \textbf{회귀 잔차}: 전문가(+0.46$\sim$+0.73)는 양수, 관리자는 0 이상, 농림어업($-2.7\sim-3.6$)과 서비스는 음수다. 사무는 노출도에 따라 다르다. Eloundou $\beta$ 기준으로는 $-0.72$(예측치의 약 절반)지만, Observed exposure와 Felten 기준으로는 0에 가깝다. Observed exposure는 AEI 사용 자료가 일부 들어간 지표라 사무직의 낮은 사용이 이미 반영되어 있다.
\end{itemize}

\section{해석상 주의}

\begin{enumerate}[leftmargin=*]
  \item \textbf{격차는 측정 오차를 포함한다.} 점유율과 노출도 모두 SOC 값을 연계표로 옮긴 것이고, 옮기는 방식이 다르다(점유율은 $w^{\text{alloc}}$ 배분, 노출도는 $w^{\text{mean}}$ 평균). 213, 293, 852, 799처럼 격차가 큰 직업 일부는 연계 경로에서 격차가 생겼을 수 있으므로, 해석 전에 어떤 SOC가 기여했는지 확인해야 한다.
  \item \textbf{AEI의 직업 분류.} AEI는 대화를 사용자 직업이 아닌 과업이 속한 SOC로 분류한다. 사무직 대화가 전문가 SOC로 분류되는 경로(paper 28)가 있으면, 사무직의 음의 격차는 과장된다.
  \item \textbf{Observed exposure는 독립적이지 않다.} AEI 사용 자료가 일부 들어 있어 점유율과의 격차가 작게 나오는 것이 자연스럽다. 사전적 노출도와의 격차는 Felten, Eloundou $\beta$ 쪽이 더 깔끔하다.
  \item \textbf{방법 선택.} ``어느 직업이 노출도에 비해 AI를 많이$\cdot$적게 쓰는가''(직업의 성격)를 보려면 회귀 잔차와 순위 격차를, ``AI 영향을 실제 사용으로 배분하면 어느 직업의 몫이 얼마나 바뀌는가''(일자리 수 시나리오)를 보려면 점유율 차이를 쓴다.
  \item \textbf{인과 해석 불가.} 격차는 한 시점(2026년 4$\cdot$5월 사용, 2022년 취업자) 단면의 차이다. 이 격차가 이후 고용 변화와 관련되는지는 별도 분석이 필요하다.
\end{enumerate}

\section{산출 파일}

\begin{itemize}[leftmargin=*]
  \item \texttt{code/04\_analysis\_kor\_share\_exposure\_gap.do}: 격차 계산, 회귀, 일치도, 목록을 만든다.
  \item \texttt{results/table/29\_kor\_share\_exposure\_gap.xlsx}: 시트 \texttt{ksco3}(129개 소분류의 지표, 순위, 격차 30개), \texttt{reg}, \texttt{agree}(모든 지표의 방법 쌍 상관), \texttt{abc}, \texttt{top10}(지표 A$\sim$C $\times$ 노출도 $\times$ 방법(로그 비율 포함)의 상위$\cdot$하위 10, 긴 형식), \texttt{robust}(지표 C 목록 8개에 든 횟수, 129개 전부), \texttt{ksco1}.
  \item \texttt{results/table/29\_reg.tex}, \texttt{29\_size.tex}, \texttt{29\_agree.tex}, \texttt{29\_abc.tex}, \texttt{29\_top\_\{OE,FEL,ELO\}.tex}, \texttt{29\_robust\_\{pos,neg\}.tex}, \texttt{29\_ksco1.tex}: 표~\ref{tab:reg}--\ref{tab:ksco1}.
\end{itemize}

\end{document}
